%!TEX program = xelatex
\let\stop\empty

\documentclass[12pt,a4paper]{exam-zh}

\usepackage{xeCJK}
\setCJKmainfont{Noto Serif CJK SC}
\usepackage{geometry}
\usepackage{array}
\usepackage{tikz}
\usepackage{amsmath}
\usepackage{wrapfig}
\usepackage{adjustbox}
\usepackage{enumitem}
\usepackage{fancyhdr}
\usepackage{pgfplots}
\pgfplotsset{compat=1.18}
\usetikzlibrary{calc}

\usepackage{scalerel}
\def\mcirc{\mathbin{\scalerel*{\circ}{j}}}
\def\msquare{\mathord{\scalerel*{\Box}{gX}}}
\geometry{margin=2cm}

\AtBeginDocument{\zihao{4}}
% ----------- 调整 exam -----------


%region exam setup
\examsetup{
    % problem/label = \arabic*., 
    % 设置大题题号格式 
        question = {
        top-sep = 1em,      % 题头上方的间距（默认 0.25em）
        bottom-sep = 1em,   % 题号下方的间距（默认 0pt）
    },
    % 隐藏大题后的分值显示 
    problem/show-points = false,
    % 调整小题列表的缩进和间距
    list/step-name = {}, 
    solution = {
    show-solution = show-move, % 答案自动移动到卷末
    label-content = {answer},  % 答案区的标题文字
    show-qed = true            % 在解答末尾显示结束符号 ◻ 
  },
    page = {
    size = a4paper,           % 设定纸张大小 
    foot-content = { - ; - }  % 页脚显示为 -1-, -2- 格式 
  }
}
\setlist[enumerate,1]{
    label=(\alpha*), 
    itemsep=0.5em, 
    leftmargin=2em, 
    after=\hfill % 配合您之前要求的百分比靠右样式
}
\newlist{questionparts}{enumerate}{1}
\setlist[questionparts]{
    label=(\alph*),
    itemsep=0.5em,
    leftmargin=2em,
    rightmargin=4em
}
\newcommand{\qscore}[1]{%
    \unskip\nobreak\hfill
    \makebox[0pt][l]{\makebox[4em][r]{(#1\%)}}%
}
\newcommand{\solutionlayout}{%
    \setlength{\abovedisplayskip}{0.6em}%
    \setlength{\belowdisplayskip}{0.6em}%
    \setlength{\abovedisplayshortskip}{0.35em}%
    \setlength{\belowdisplayshortskip}{0.35em}%
    \setlength{\jot}{0.35em}%
}
%endregion




\begin{document}

%region 范围：普通高中数学 | 考点：数学归纳法证明恒等式（求和公式） | 难度：基础
\begin{question}[points = 8]
    用数学归纳法证明：对任意正整数 $n$，
    \[
    1+2+3+\cdots+n=\dfrac{n(n+1)}{2}.
    \]
    \hfill (8\%)
\end{question}

\begin{solution}
    证明对任意正整数 $n$，$1+2+\cdots+n=\dfrac{n(n+1)}{2}$ 成立。\\
    按数学归纳法，先验证 $n=1$ 时等式成立，再证明由 $n=k$ 成立可推出 $n=k+1$ 成立。\\
    当 $n=1$ 时，
    \[
    \text{左边}=1,\qquad \text{右边}=\dfrac{1\cdot 2}{2}=1.
    \]
    \hfill (2分) \\
    等式成立。假设 $n=k$（$k\geqslant 1$）时命题成立，即
    \[
    1+2+\cdots+k=\dfrac{k(k+1)}{2}.
    \]
    \hfill (2分) \\
    则当 $n=k+1$ 时，
    \begin{align*}
        1+2+\cdots+k+(k+1) &= \dfrac{k(k+1)}{2}+(k+1) \\
        &= \dfrac{k(k+1)+2(k+1)}{2} \\
        &= \dfrac{(k+1)(k+2)}{2}.
    \end{align*}
    \hfill (3分) \\
    这恰为 $n=k+1$ 时等式右边 $\dfrac{(k+1)\left[(k+1)+1\right]}{2}$ 的形式，故 $n=k+1$ 时命题也成立。由数学归纳法原理，对任意正整数 $n$，原等式成立。
    \hfill (1分)
\end{solution}
%endregion

%region 范围：普通高中数学 | 考点：数学归纳法证明整除性 | 难度：基础
\begin{question}[points = 8]
    用数学归纳法证明：对任意正整数 $n$，$5^n-1$ 能被 $4$ 整除。
    \hfill (8\%)
\end{question}

\begin{solution}
    证明对任意正整数 $n$，$4\mid(5^n-1)$ 成立。\\
    按数学归纳法，先验证 $n=1$ 时命题成立，再证明由 $n=k$ 成立可推出 $n=k+1$ 成立。\\
    当 $n=1$ 时，
    \[
    5^1-1=4,
    \]
    显然能被 $4$ 整除。
    \hfill (2分) \\
    假设 $n=k$（$k\geqslant 1$）时命题成立，即存在正整数 $m$，使得
    \[
    5^k-1=4m.
    \]
    \hfill (2分) \\
    则当 $n=k+1$ 时，
    \begin{align*}
        5^{k+1}-1 &= 5\cdot 5^k-1 \\
        &= 5(4m+1)-1 \\
        &= 20m+4 \\
        &= 4(5m+1).
    \end{align*}
    \hfill (3分) \\
    因 $5m+1$ 为正整数，故 $5^{k+1}-1$ 能被 $4$ 整除，即 $n=k+1$ 时命题也成立。由数学归纳法原理，对任意正整数 $n$，$5^n-1$ 能被 $4$ 整除。
    \hfill (1分)
\end{solution}
%endregion

%region 范围：普通高中数学 | 考点：数学归纳法证明恒等式 | 难度：中等
\begin{question}[points = 10]
    用数学归纳法证明：对任意正整数 $n$，
    \[
    1^2+2^2+3^2+\cdots+n^2=\dfrac{n(n+1)(2n+1)}{6}.
    \]
    \hfill (10\%)
\end{question}

\begin{solution}
    证明对任意正整数 $n$，$1^2+2^2+\cdots+n^2=\dfrac{n(n+1)(2n+1)}{6}$ 成立。\\
    按数学归纳法，先验证 $n=1$ 时等式成立，再证明由 $n=k$ 成立可推出 $n=k+1$ 成立。\\
    当 $n=1$ 时，
    \[
    \text{左边}=1^2=1,\qquad \text{右边}=\dfrac{1\cdot 2\cdot 3}{6}=1.
    \]
    \hfill (2分) \\
    等式成立。假设 $n=k$（$k\geqslant 1$）时命题成立，即
    \[
    1^2+2^2+\cdots+k^2=\dfrac{k(k+1)(2k+1)}{6}.
    \]
    \hfill (2分) \\
    则当 $n=k+1$ 时，
    \begin{align*}
        1^2+2^2+\cdots+k^2+(k+1)^2 &= \dfrac{k(k+1)(2k+1)}{6}+(k+1)^2 \\
        &= \dfrac{k(k+1)(2k+1)+6(k+1)^2}{6} \\
        &= \dfrac{(k+1)\left[k(2k+1)+6(k+1)\right]}{6} \\
        &= \dfrac{(k+1)(2k^2+7k+6)}{6} \\
        &= \dfrac{(k+1)(k+2)(2k+3)}{6}.
    \end{align*}
    \hfill (4分) \\
    这恰为 $n=k+1$ 时等式右边 $\dfrac{(k+1)\left[(k+1)+1\right]\left[2(k+1)+1\right]}{6}$ 的形式，故 $n=k+1$ 时命题也成立。由数学归纳法原理，对任意正整数 $n$，原等式成立。
    \hfill (2分)
\end{solution}
%endregion

%region 范围：普通高中数学 | 考点：数学归纳法证明不等式 | 难度：中等
\begin{question}[points = 10]
    用数学归纳法证明：当 $n\geqslant 5$ 时，
    \[
    2^n>n^2.
    \]
    \hfill (10\%)
\end{question}

\begin{solution}
    证明当 $n\geqslant 5$ 时 $2^n>n^2$ 恒成立。\\
    按数学归纳法，先验证 $n=5$ 时不等式成立，再证明由 $n=k$ 成立可推出 $n=k+1$ 成立。\\
    当 $n=5$ 时，
    \[
    2^5=32,\qquad 5^2=25,
    \]
    因 $32>25$，不等式成立。
    \hfill (2分) \\
    假设 $n=k$（$k\geqslant 5$）时命题成立，即
    \[
    2^k>k^2.
    \]
    \hfill (2分) \\
    则当 $n=k+1$ 时，
    \begin{align*}
        2^{k+1} &= 2\cdot 2^k \\
        &> 2k^2.
    \end{align*}
    \hfill (2分) \\
    又因 $k\geqslant 5$，有
    \begin{align*}
        2k^2-(k+1)^2 &= k^2-2k-1 \\
        &= (k-1)^2-2 \\
        &>0,
    \end{align*}
    故 $2k^2>(k+1)^2$，从而 $2^{k+1}>(k+1)^2$，即 $n=k+1$ 时命题也成立。由数学归纳法原理，当 $n\geqslant 5$ 时，$2^n>n^2$ 恒成立。
    \hfill (4分)
\end{solution}
%endregion

%region 范围：普通高中数学 | 考点：数学归纳法证明不等式（根式型） | 难度：困难
\begin{question}[points = 12]
    用数学归纳法证明：当 $n\geqslant 2$ 时，
    \[
    1+\dfrac{1}{\sqrt{2}}+\dfrac{1}{\sqrt{3}}+\cdots+\dfrac{1}{\sqrt{n}}>\sqrt{n}.
    \]
    \hfill (12\%)
\end{question}

\begin{solution}
    证明当 $n\geqslant 2$ 时 $\displaystyle 1+\frac{1}{\sqrt2}+\cdots+\frac{1}{\sqrt n}>\sqrt n$ 恒成立。\\
    按数学归纳法，先验证 $n=2$ 时不等式成立，再利用 $\sqrt{k+1}-\sqrt{k}=\dfrac{1}{\sqrt{k+1}+\sqrt{k}}$ 完成归纳步。\\
    当 $n=2$ 时，
    \begin{align*}
        1+\dfrac{1}{\sqrt2}-\sqrt2 &= \dfrac{\sqrt2+1-2}{\sqrt2} \\
        &= \dfrac{\sqrt2-1}{\sqrt2} \\
        &>0,
    \end{align*}
    故 $1+\dfrac{1}{\sqrt2}>\sqrt2$，不等式成立。
    \hfill (3分) \\
    假设 $n=k$（$k\geqslant 2$）时命题成立，即
    \[
    1+\dfrac{1}{\sqrt2}+\cdots+\dfrac{1}{\sqrt k}>\sqrt k.
    \]
    \hfill (2分) \\
    则当 $n=k+1$ 时，
    \[
    1+\dfrac{1}{\sqrt2}+\cdots+\dfrac{1}{\sqrt k}+\dfrac{1}{\sqrt{k+1}}>\sqrt k+\dfrac{1}{\sqrt{k+1}}.
    \]
    \hfill (2分) \\
    只需证明 $\sqrt k+\dfrac{1}{\sqrt{k+1}}\geqslant \sqrt{k+1}$，即证
    \begin{align*}
        \dfrac{1}{\sqrt{k+1}} &\geqslant \sqrt{k+1}-\sqrt k \\
        &= \dfrac{1}{\sqrt{k+1}+\sqrt k}.
    \end{align*}
    \hfill (3分) \\
    因 $k\geqslant 2$，有 $\sqrt{k+1}+\sqrt k>\sqrt{k+1}$，故 $\dfrac{1}{\sqrt{k+1}+\sqrt k}<\dfrac{1}{\sqrt{k+1}}$，上式成立。从而
    \[
    1+\dfrac{1}{\sqrt2}+\cdots+\dfrac{1}{\sqrt{k+1}}>\sqrt{k+1},
    \]
    即 $n=k+1$ 时命题也成立。由数学归纳法原理，当 $n\geqslant 2$ 时，原不等式恒成立。
    \hfill (2分)
\end{solution}
%endregion

\newpage
\thispagestyle{empty}
\null
\newpage

\end{document}